\documentclass[11pt]{article}

\usepackage[a4paper,margin=1in]{geometry}
\usepackage{amsmath,amssymb,amsthm}
\usepackage{bm}
\usepackage{booktabs}
\usepackage{xcolor}
\usepackage{microtype}
\usepackage[colorlinks=true,linkcolor=teal!60!black,citecolor=teal!60!black,urlcolor=teal!60!black]{hyperref}
\usepackage{enumitem}
\usepackage{titlesec}

\definecolor{quasarblue}{HTML}{0B5563}
\definecolor{verified}{HTML}{1B7A43}
\definecolor{phenom}{HTML}{B26B00}
\definecolor{architect}{HTML}{6A4C93}

\titleformat{\section}{\large\bfseries\color{quasarblue}}{\thesection}{0.6em}{}
\titleformat{\subsection}{\normalsize\bfseries\color{quasarblue!85!black}}{\thesubsection}{0.6em}{}

% Status tags that distinguish literature-grounded equations from Quasar's own modelling
\newcommand{\statusVerified}{\textbf{\color{verified}[\,LITERATURE-GROUNDED\,]}}
\newcommand{\statusPhenom}{\textbf{\color{phenom}[\,PHENOMENOLOGICAL MODEL\,]}}
\newcommand{\statusArch}{\textbf{\color{architect}[\,ARCHITECTURAL COUPLING\,]}}

\newcommand{\Tr}{\operatorname{Tr}}
\newcommand{\dd}{\mathrm{d}}
\newcommand{\ii}{\mathrm{i}}

\title{\vspace{-1.2cm}\textbf{\color{quasarblue}The Physics of Quasar}\\[2pt]
\large A Quantum Network Digital Twin: Governing Equations and Their Sources}
\author{Technical Reference --- Quasar (\texttt{qndt})}
\date{\today}

\begin{document}
\maketitle
\vspace{-0.7cm}

\begin{abstract}
\noindent Quasar is a Quantum Network Digital Twin (QNDT) that models a fielded quantum
communication link not as an idealised textbook channel but as a time-varying,
telemetry-driven cyber-physical asset. This document sets out the mathematical model
the simulator implements, module by module, and traces each governing equation to its
source in the open-systems and quantum-communication literature. Throughout, a
deliberate distinction is drawn between equations that reproduce established, peer-reviewed
results and equations that are Quasar's own \emph{phenomenological} engineering models.
The latter are physically motivated but are not, and should not be read as, derivations
from a single canonical paper. Each module is tagged accordingly. The intent is a
faithful account of what the simulator computes and why --- not a claim that every line
of the model is settled physics.
\end{abstract}

\noindent\textbf{Reading the status tags.}\quad
\statusVerified{} the equation reproduces a result with a cited primary source.\;
\statusPhenom{} a physically-motivated model internal to Quasar, calibrated rather than derived.\;
\statusArch{} a coupling rule defining how subsystems feed one another in the twin.

\section{Scope and design premise}
Discrete-event quantum-network packages such as NetSquid, SeQUeNCe and QuISP typically
parameterise the channel as a quasi-stationary object, even where some time-dependence can be
configured. Quasar's premise is to make non-stationarity primary: the channel is a \emph{stochastic, non-stationary} process whose instantaneous
error structure is driven by live environmental telemetry (seismic, thermal, wind-induced
polarisation drift), by classical co-propagating traffic, and by the operational history of
the hardware. The model therefore rests on five pillars: (i) a Pauli-channel representation
of single-link noise; (ii) a telemetry-driven, non-Markovian time-varying channel;
(iii) a classical$\,\to\,$quantum coexistence (Raman) noise engine; (iv) stateful device aging;
and (v) a tensor-network reduced-order state backend. The sections below treat each in turn.

\section{Single-link noise: the Pauli-channel and its transfer matrix}
\statusVerified

A quantum channel acting on a density operator $\rho$ is a completely-positive,
trace-preserving (CPTP) map, expressible in Kraus form \cite{nielsenchuang}
\begin{equation}
\Phi(\rho)=\sum_i A_i\,\rho\,A_i^{\dagger},
\qquad \sum_i A_i^{\dagger}A_i=\mathbb{1}.
\end{equation}
Quasar represents each link's noise in the \emph{Pauli Transfer Matrix} (PTM) basis, which
expresses $\Phi$ as a $d^2\!\times\!d^2$ real matrix in the Pauli basis
$\{P_j\}$ (with $P_0=\mathbb{1}$ and the $1/d$ prefactor absorbing the basis normalisation),
$d=2^n$ \cite{ptm_recon}:
\begin{equation}
\big(R_{\Phi}\big)_{ij}=\frac{1}{d}\,\Tr\!\big[P_i\,\Phi(P_j)\big],
\qquad |\Phi(\rho)\rangle\!\rangle = R_{\Phi}\,|\rho\rangle\!\rangle,
\end{equation}
where $|\cdot\rangle\!\rangle$ denotes the column-vectorised operator. The PTM is preferred
because channel \emph{composition becomes matrix multiplication}, and the experimentally
accessible quantities (visibility / contrast decay) appear directly as its eigenvalues.

For a single-qubit \emph{Pauli channel} --- the workhorse class spanning the depolarising,
dephasing and bit-flip models --- the map and its eigenrelation read \cite{stt_pauli}
\begin{equation}
\Lambda(\rho)=\!\!\sum_{\alpha\in\{0,x,y,z\}}\!\! f_\alpha\,P_\alpha\,\rho\,P_\alpha,
\qquad \sum_\alpha f_\alpha = 1,
\qquad \Lambda(P_\alpha)=\lambda_\alpha P_\alpha,\;\; \lambda_0=1 .
\end{equation}
The probabilities $f_\alpha$ and the PTM eigenvalues $\lambda_\alpha$ are related by a
Walsh--Hadamard transform,
\begin{equation}
\lambda_x=f_0+f_x-f_y-f_z,\quad
\lambda_y=f_0-f_x+f_y-f_z,\quad
\lambda_z=f_0-f_x-f_y+f_z .
\end{equation}
Hence the PTM of any Pauli channel is \emph{diagonal},
\begin{equation}
\boxed{\;R=\operatorname{diag}(1,\lambda_x,\lambda_y,\lambda_z)\;}
\end{equation}
and the composition of two Pauli channels reduces to an element-wise (Hadamard) product of
their eigenvalues,
\begin{equation}
\lambda^{(2\circ1)}_\alpha=\lambda^{(2)}_\alpha\,\lambda^{(1)}_\alpha ,
\end{equation}
an $O(1)$, numerically exact operation per Pauli component. This is the algebra Quasar uses
to compose the noise contributions of consecutive fibre segments and devices along a path.

\section{Telemetry-driven, time-varying channel}
\statusPhenom{} (structure) / \statusVerified{} (memory-kernel formalism)

The defining novelty of Quasar is that the Pauli rates are not constants but functionals of
an environmental state vector $\bm{E}(t)$ (e.g.\ temperature, strain, wind-loading proxies).
The instantaneous Pauli-rate vector $\bm{p}(t)=(p_x,p_y,p_z)$ is obtained by convolving the
environmental drive against a memory kernel and passing the result through a link function
$\Phi_{\mathrm{link}}$:
\begin{equation}
\boxed{\;\bm{p}(t)=\Phi_{\mathrm{link}}\!\left(\int_{-\infty}^{t}
\mathbf{K}(t-t')\,\mathbf{S}\,\bm{E}(t')\,\dd t'\right).\;}
\end{equation}
Here $\mathbf{S}$ is a \emph{sensitivity matrix} mapping physical environmental channels onto
Pauli-stress axes, and $\mathbf{K}$ is a matrix of memory kernels. The structure of this
equation --- and the matrix $\mathbf{S}$ in particular --- is Quasar's own modelling choice.
The shipped default values of $\mathbf{S}$ are physically-motivated \emph{illustrative} figures,
not measurements from a real fibre link: no field calibration data, literature citation, or
derivation underpins their magnitudes.  They are chosen to produce a legible demonstration
QBER at $20\,^\circ\mathrm{C}$ Celsius with the normalised kernel; note that the temperature
channel couples to the \emph{absolute} Celsius value, not to a $\Delta T$, so the resulting
QBER depends on the temperature-scale origin rather than on a physical drift.  The sensitivity
structure is tagged phenomenological for that reason.
What \emph{is} literature-grounded is the kernel formalism: a present-time response that depends
on the integrated past is the hallmark of \emph{non-Markovian} dynamics
\cite{breuerbook,rhp2014}. Quasar ships three standard kernel shapes, each multiplied by a
normalisation constant fixed so that $\int_0^{\infty}K(\tau)\,\dd\tau=1$:
\begin{align}
\text{Exponential:}\quad & K(\tau)=\frac{1}{\tau_c}\,e^{-\tau/\tau_c}\,\Theta(\tau),\\
\text{Damped oscillatory (Lorentzian spectrum):}\quad & K(\tau)=\mathcal{N}_{\!L}\,e^{-\gamma\tau}\cos(\Omega\tau)\,\Theta(\tau),\\
\text{Gaussian (one-sided):}\quad & K(\tau)=\mathcal{N}_{\!G}\,e^{-\tau^2/2\sigma^2}\,\Theta(\tau),
\end{align}
with $\Theta$ the Heaviside step enforcing causality and
$\mathcal{N}_{\!L}=(\gamma^2+\Omega^2)/\gamma$, $\mathcal{N}_{\!G}=\sqrt{2/\pi\sigma^2}$ the
half-line normalisation constants. (The exponential kernel is already unit-area, so its
constant is unity.) The second kernel is named for its \emph{Lorentzian spectral density}, not
its time-domain shape; the time-domain form is a damped oscillation. The exponential kernel
recovers the memoryless (Markovian) limit as $\tau_c\!\to\!0$; the damped-oscillatory kernel
admits information back-flow through its sign changes; the Gaussian models a finite,
smoothly-windowed memory. The precise parameterisation and normalisation convention should be
matched against the shipped \texttt{kernels.py} rather than assumed from these forms.

\section{Open-system dynamics: TCL master equation and the non-Markovianity witness}
\statusVerified

When a full density-matrix evolution is required rather than a static Pauli map, Quasar
integrates a time-local (time-convolutionless, TCL) master equation in
Gorini--Kossakowski--Sudarshan--Lindblad (GKSL) form with \emph{time-dependent} rates
\cite{lindblad1976,gks1976,breuerbook}:
\begin{equation}
\boxed{\;\dot\rho(t)=-\frac{\ii}{\hbar}\big[H(t),\rho(t)\big]
+\sum_k \gamma_k(t)\!\left(L_k\,\rho\,L_k^{\dagger}
-\tfrac{1}{2}\big\{L_k^{\dagger}L_k,\rho\big\}\right).\;}
\end{equation}
The decisive feature is the sign of the canonical rates $\gamma_k(t)$. The map is
CP-divisible --- i.e.\ Markovian in the Rivas--Huelga--Plenio sense --- if and only if
$\gamma_k(t)\ge 0$ for all $k$ and all $t$. \emph{Temporary negativity} of any
$\gamma_k(t)$ is the signature of non-Markovian memory and information back-flow from the
environment to the system \cite{rhp2010,rhp2014}.

\subsection{Quantifying memory: the RHP measure and Quasar's online witness}
Rivas, Huelga and Plenio define Markovian dynamics as a CP-divisible family of propagators
obeying the composition law $\mathcal{E}_{(t_2,t_0)}=\mathcal{E}_{(t_2,t_1)}\mathcal{E}_{(t_1,t_0)}$,
which holds \emph{if and only if} the time-local generator can be written in the standard form
above with all canonical rates $\gamma_k(t)\ge 0$ \cite{rhp2010}. Their tomography-free measure
quantifies the resulting non-monotonicity of system--ancilla entanglement. For a system
initially maximally entangled with a shielded ancilla, $\rho_{SA}(t_0)=|\Phi\rangle\!\langle\Phi|$,
and an entanglement measure $E$, it reads \cite{rhp2010}
\begin{equation}
\mathcal{I}^{(E)}=\int_{t_0}^{t_{\max}}\left|\frac{\dd E[\rho_{SA}(t)]}{\dd t}\right|\dd t
-\Delta E,
\qquad \Delta E = E[\rho_{SA}(t_0)]-E[\rho_{SA}(t_{\max})],
\end{equation}
which vanishes when the entanglement decays monotonically (Markovian case) and is strictly
positive whenever entanglement is transiently regained.

Quasar's generator is constructed directly in time-local form, so it has access to the
canonical rates themselves and does not perform ancilla tomography. It therefore reports a
\emph{rate-based divisibility witness},
\begin{equation}
\boxed{\;\mathcal{N}_{\mathrm{rate}}=\sum_k\int_{\gamma_k(t)<0}\big|\gamma_k(t)\big|\,\dd t,
\qquad \mathcal{N}_{\mathrm{rate}}>0 \;\Longleftrightarrow\; \exists\,k,t:\;\gamma_k(t)<0,\;}
\end{equation}
i.e.\ a violation of CP-divisibility and hence of Markovianity in the RHP sense. This is the
quantity the simulator displays. It must be read for what it is: an internal Quasar diagnostic
built on the RHP divisibility \emph{criterion}, not a numerical reconstruction of
$\mathcal{I}^{(E)}$, and its absolute magnitude is not on the same scale as the
entanglement-based measure. A trace-distance diagnostic following Breuer--Laine--Piilo
\cite{blp2009} is also available where only state-level (non-tomographic) data exist. Fast
thermal transients and seismic events are the physical triggers that drive $\gamma_k(t)$
negative in the model.

\section{Classical$\,\to\,$quantum coexistence: the Raman interference engine}
\statusVerified

When the quantum channel shares fibre with classical WDM traffic, spontaneous Raman
scattering (SpRS) of the strong classical pump generates broadband in-band photons that are
spectrally indistinguishable from the single-photon signal and therefore \emph{cannot be
filtered out} \cite{townsend1997,eraerds2010}. Quasar models the Raman power collected in the
quantum-channel detection bandwidth $\Delta\lambda$ for the two propagation geometries as
\cite{dasilva2014,peters2009}
\begin{align}
\text{Forward (co-propagating):}\quad & P_{\mathrm{f}}(L)=P_0\,C_{\mathrm{R}}(\Delta\nu)\,\Delta\lambda\,L\,e^{-\alpha L},\\
\text{Backward (counter-propagating):}\quad & P_{\mathrm{b}}(L)=P_0\,C_{\mathrm{R}}(\Delta\nu)\,\Delta\lambda\,\frac{1-e^{-2\alpha L}}{2\alpha},
\end{align}
where $P_0$ is the launched classical power, $L$ the fibre length, $\alpha$ the attenuation
coefficient (taken equal at the pump and scattered wavelengths --- the standard approximation),
and $C_{\mathrm{R}}(\Delta\nu)$ the SpRS cross-section per unit length per unit bandwidth
at frequency offset $\Delta\nu$ between the classical and quantum channels. Near the quantum
channel $C_{\mathrm{R}}$ is well approximated by a V-shaped linear law with distinct Stokes /
anti-Stokes slopes \cite{dasilva2014},
\begin{equation}
C_{\mathrm{R}}(\Delta\nu)\simeq
\begin{cases}
C_0+\beta_{\mathrm{S}}\,|\Delta\nu|, & \nu_{\mathrm{cl}}>\nu_{\mathrm{q}} \;(\text{Stokes}),\\[2pt]
C_0+\beta_{\mathrm{AS}}\,|\Delta\nu|, & \nu_{\mathrm{cl}}<\nu_{\mathrm{q}} \;(\text{anti-Stokes}).
\end{cases}
\end{equation}
For $M$ active classical channels the contributions add incoherently,
$P_{\mathrm{R}}=\sum_{m=1}^{M}P^{(m)}$. The collected Raman power is then converted into a mean
induced detector count per gate of length $\Delta t$ at detection efficiency $\eta$, with
$h\nu$ the quantum-channel photon energy (this mean equals a per-gate click probability in the
relevant regime $p_{\mathrm{R}}\ll 1$),
\begin{equation}
\boxed{\;p_{\mathrm{R}}=\frac{\eta\,\Delta t}{h\nu}\,P_{\mathrm{R}},\qquad
Y_0 = 2\,p_{\mathrm{dark}}+p_{\mathrm{R}},\;}
\end{equation}
so that classical traffic raises the effective background floor $Y_0$ that feeds directly into
the link's QBER and secure-key-rate estimates \cite{eraerds2010}. Quasar models a
BB84-style dual-detector receiver: the two single-photon detectors contribute independent
intrinsic dark counts, hence the $2\,p_{\mathrm{dark}}$ multiplicity, whereas the Raman term
$p_{\mathrm{R}}$ enters once, since an incident Raman photon is a single photon registering on
one detector arm. (A single-detector or single-arm receiver would instead use
$Y_0=p_{\mathrm{dark}}+p_{\mathrm{R}}$.) This is the mechanism by which a busy classical fibre
silently degrades a co-existing quantum channel.

\section{Stateful device aging and imperfect calibration}
\statusPhenom

Memory nodes in Quasar carry a coherence budget that \emph{erodes with use}. Idealised
dephasing imposes the transverse Pauli eigenvalues $\lambda_x(t)=\lambda_y(t)=e^{-t/T_2}$
(the longitudinal eigenvalue $\lambda_z$ is governed by $T_1$); Quasar makes $T_2$ itself a
function of accumulated operational duty cycle $D(t)=\int_0^t u(t')\,\dd t'$, where $u$ is the
instantaneous utilisation:
\begin{equation}
\boxed{\;\frac{1}{T_2\!\big(D\big)}=\frac{1}{T_2^{(0)}}+\kappa\,D(t),\qquad
\frac{1}{T_2}=\frac{1}{2T_1}+\frac{1}{T_\phi}.\;}
\end{equation}
The second identity is the standard relation linking the coherence time $T_2$, the
relaxation time $T_1$ and the pure-dephasing time $T_\phi$ \cite{nielsenchuang}; it is exact
and is included so the wear law has an unambiguous effect on the dephasing channel. The
\emph{form} of the dephasing eigenvalue and this $T_1/T_2/T_\phi$ decomposition are standard.
The duty-cycle wear term $\kappa\,D$, by contrast, is Quasar's own engineering heuristic ---
calibrated against device logs, not derived --- which is why the module as a whole carries the
phenomenological tag. Imperfect calibration is modelled by sampling the as-built parameters
$(T_2^{(0)},T_1,\eta,\dots)$ from device-specific tolerance distributions rather than fixing
them at nominal values.

\section{Classical--quantum control-plane desynchronisation}
\statusArch

The control plane runs as an asynchronous discrete-event process. Congestion in the classical
signalling layer --- routing churn, packet jitter, queueing delay $W(t)$ --- lengthens the
time a quantum memory must hold its state before a gate or swap can be scheduled. Quasar
couples the two planes through an effective hold time,
\begin{equation}
\tau_{\mathrm{hold}}(t)=\tau_{\mathrm{base}}+W(t)\quad\Longrightarrow\quad
\lambda_x,\lambda_y=\exp\!\big(-\tau_{\mathrm{hold}}(t)/T_2(t)\big),
\end{equation}
so that classical-network pathology propagates into quantum-state degradation. This is a
structural coupling within the twin rather than a fundamental physical law; it is included to
capture a failure mode that stationary-channel simulators cannot express.

\section{Reduced-order state backend: matrix-product density operators}
\statusVerified

Tracking a global density matrix costs $O(d^{2N})$ for $N$ nodes and is intractable at network
scale. Quasar's state backend instead represents the joint density operator as a
\emph{matrix-product density operator} (MPDO)
\cite{schollwock2011,weimer2021},
\begin{equation}
\rho=\sum_{\{i_n,j_n\}} M_1^{i_1 j_1}\,M_2^{i_2 j_2}\cdots M_N^{i_N j_N}\;
|i_1\cdots i_N\rangle\langle j_1\cdots j_N|,
\end{equation}
with bond dimension $\chi$ controlling the retained correlations. Singular-value truncation at
each bond keeps $\chi$ bounded, reducing the cost from exponential to
$O\!\big(N\chi^3 d^2\big)$ while preserving the trace and Hermiticity of $\rho$
\cite{jarkovsky2020}. Positive-semidefiniteness, by contrast, is guaranteed by construction
only in the locally-purified variant (LPTN); plain MPDO truncation can in principle introduce
small negative eigenvalues, so Quasar monitors the minimum eigenvalue as a validity check. The
exact prefactor and power of $\chi$ are operation-dependent; the
load-bearing point is that the cost grows linearly in $N$ and polynomially (rather than
exponentially) in the retained bond dimension. Accuracy is therefore tunable: $\chi$ trades
fidelity against speed, and the truncation error is the discarded singular weight, which
Quasar can report alongside each evolution step.

\section{Validation status and epistemic scope}
This document describes the model Quasar is \emph{designed} to implement; it is not a
certificate that the shipped code is free of deviation from these equations. Two honesty
caveats follow directly from that distinction. First, the literature-grounded modules
(Pauli-channel algebra, the GKSL/TCL master equation, the RHP non-Markovianity measure, the
SpRS power laws, the MPDO backend) can be checked line-by-line against their primary sources
and against analytic limiting cases --- and should be. Second, the phenomenological modules (the
sensitivity matrix $\mathbf{S}$, the duty-cycle wear law, the control-plane coupling) are
physically motivated and internally consistent, but their default numerical values are
\emph{illustrative} rather than empirically calibrated: no field telemetry, literature
measurement, or formal derivation underpins the specific numbers in $\mathbf{S}$.
A rigorous audit therefore verifies the first class against citations and derives the
second class from measured fibre data, reporting any gap explicitly rather than asserting
global correctness.

Every reference in the bibliography below has been verified against the publisher's record
(APS, IOP, Springer, AIP and equivalent) for authorship, venue, volume and year; several were
checked against the primary text directly. Independently, the load-bearing algebra of this
document --- the diagonality of the Pauli-channel transfer matrix and its eigenvalues, the
dephasing eigenvalue assignment ($\lambda_x=\lambda_y=e^{-t/T_2}$ with $\lambda_z$ governed by
$T_1$), channel composition as a Hadamard product, the forward/backward Raman power integrals,
and the kernel normalisation constants --- has been re-derived by independent symbolic
computation and confirmed (Appendix~A). Where an equation is standard but its exact
normalisation or sign convention is implementation-dependent --- the memory-kernel
normalisation and the dephasing-rate bookkeeping --- the text says so and defers to the shipped
code as ground truth. The receiver model is fixed by design as a BB84-style dual-detector
configuration ($Y_0=2\,p_{\mathrm{dark}}+p_{\mathrm{R}}$); the shipped code must match this
rather than the single-detector form. No equation in this document should be read as a
claim about the running implementation that has not itself been checked against the
corresponding module.

\appendix
\section{Independent symbolic verification}
The identities below were checked by computer algebra (\texttt{sympy}/\texttt{numpy}); all
passed exactly (not merely numerically to tolerance, except where a numerical example is noted).

\begin{enumerate}[itemsep=2pt,topsep=2pt,leftmargin=1.5em]
\item \textbf{PTM diagonality.} For a generic single-qubit Pauli channel
$\Lambda(\rho)=\sum_\alpha f_\alpha P_\alpha\rho P_\alpha$, the matrix
$(R)_{ij}=\tfrac{1}{d}\Tr[P_i\Lambda(P_j)]$ is diagonal, $R=\mathrm{diag}(1,\lambda_x,\lambda_y,\lambda_z)$.
\item \textbf{Walsh--Hadamard eigenvalues.} The formulas
$\lambda_x=f_0+f_x-f_y-f_z$ (and cyclic) reproduce the computed PTM diagonal exactly.
\item \textbf{Dephasing assignment.} The phase-flip channel ($f_0=1-p,\,f_z=p$) yields
$\lambda_x=\lambda_y=1-2p$ (the transverse coherence, $\to e^{-t/T_2}$) and $\lambda_z=1$,
confirming that $T_2$ decay attaches to $\lambda_x,\lambda_y$ --- not $\lambda_z$.
\item \textbf{Composition.} For two Pauli channels, $\mathrm{diag}(R_2R_1)=\mathrm{diag}(R_2)\odot\mathrm{diag}(R_1)$
(Hadamard product).
\item \textbf{Raman integrals.} The distributed-scattering integrals reduce to the document's
closed forms,
\begin{align*}
P_{\mathrm{f}}(L) &= P_0\,C\,\Delta\lambda\!\int_0^L e^{-\alpha z}\,e^{-\alpha(L-z)}\,\dd z = P_0\,C\,\Delta\lambda\,L\,e^{-\alpha L},\\
P_{\mathrm{b}}(L) &= P_0\,C\,\Delta\lambda\!\int_0^L e^{-2\alpha z}\,\dd z
= P_0\,C\,\Delta\lambda\,\tfrac{1-e^{-2\alpha L}}{2\alpha}.
\end{align*}
\vspace{-0.3em}
\item \textbf{Kernel normalisations.}
\begin{align*}
\int_0^\infty e^{-\gamma\tau}\cos(\Omega\tau)\,\dd\tau=\frac{\gamma}{\gamma^2+\Omega^2}
&\;\Rightarrow\; \mathcal{N}_L=\frac{\gamma^2+\Omega^2}{\gamma},\\
\int_0^\infty e^{-\tau^2/2\sigma^2}\,\dd\tau=\sigma\sqrt{\pi/2}
&\;\Rightarrow\; \mathcal{N}_G=\sqrt{2/\pi\sigma^2}.
\end{align*}
\end{enumerate}
These checks confirm the internal algebra of the document. They do not certify the shipped
code; that requires the per-module audit described in \S9.

\begin{thebibliography}{99}
\setlength{\itemsep}{2pt}

\bibitem{nielsenchuang} M.~A. Nielsen and I.~L. Chuang,
\emph{Quantum Computation and Quantum Information}, 10th Anniversary Ed.
(Cambridge University Press, 2010).

\bibitem{ptm_recon} S.~Roncallo, L.~Maccone and C.~Macchiavello,
``Pauli transfer matrix direct reconstruction: channel characterization without process
tomography,'' \emph{Quantum Sci. Technol.} \textbf{9}, 015010 (2024); arXiv:2212.11968.
PTM/Pauli--Liouville definition $(R_\Phi)_{ij}=\tfrac{1}{d}\Tr[P_i\,\Phi(P_j)]$ used here.

\bibitem{stt_pauli} Y.-Q. Chen, Y.-C. Zheng, S.~Zhang and C.-Y. Hsieh,
``Spectral-transfer-tensor method for characterizing non-Markovian noise,''
\emph{Phys. Rev. Applied} \textbf{17}, 064007 (2022); arXiv:2012.10094.
One-qubit Pauli channel and its eigenvalue decomposition used here.

\bibitem{lindblad1976} G.~Lindblad, ``On the generators of quantum dynamical semigroups,''
\emph{Commun. Math. Phys.} \textbf{48}, 119--130 (1976).

\bibitem{gks1976} V.~Gorini, A.~Kossakowski and E.~C.~G. Sudarshan,
``Completely positive dynamical semigroups of $N$-level systems,''
\emph{J. Math. Phys.} \textbf{17}, 821--825 (1976).

\bibitem{breuerbook} H.-P. Breuer and F.~Petruccione,
\emph{The Theory of Open Quantum Systems} (Oxford University Press, 2002).

\bibitem{rhp2010} \'A.~Rivas, S.~F. Huelga and M.~B. Plenio,
``Entanglement and non-Markovianity of quantum evolutions,''
\emph{Phys. Rev. Lett.} \textbf{105}, 050403 (2010).

\bibitem{rhp2014} \'A.~Rivas, S.~F. Huelga and M.~B. Plenio,
``Quantum non-Markovianity: characterization, quantification and detection,''
\emph{Rep. Prog. Phys.} \textbf{77}, 094001 (2014).

\bibitem{blp2009} H.-P. Breuer, E.-M. Laine and J.~Piilo,
``Measure for the degree of non-Markovian behavior of quantum processes in open systems,''
\emph{Phys. Rev. Lett.} \textbf{103}, 210401 (2009).

\bibitem{townsend1997} P.~D. Townsend,
``Simultaneous quantum cryptographic key distribution and conventional data transmission over
installed fibre using wavelength-division multiplexing,''
\emph{Electron. Lett.} \textbf{33}, 188--190 (1997).

\bibitem{eraerds2010} P.~Eraerds, N.~Walenta, M.~Legr\'e, N.~Gisin and H.~Zbinden,
``Quantum key distribution and 1\,Gbps data encryption over a single fibre,''
\emph{New J. Phys.} \textbf{12}, 063027 (2010).

\bibitem{peters2009} N.~A. Peters, P.~Toliver, T.~E. Chapuran \emph{et al.},
``Dense wavelength multiplexing of 1550\,nm QKD with strong classical channels in
reconfigurable networking environments,''
\emph{New J. Phys.} \textbf{11}, 045012 (2009).

\bibitem{dasilva2014} T.~Ferreira da Silva, G.~B. Xavier, G.~P. Tempor\~ao and
J.~P. von der Weid, ``Impact of Raman scattered noise from multiple telecom channels on
fiber-optic quantum key distribution systems,''
\emph{J. Lightwave Technol.} \textbf{32}, 2332--2339 (2014).

\bibitem{schollwock2011} U.~Schollw\"ock,
``The density-matrix renormalization group in the age of matrix product states,''
\emph{Ann. Phys.} \textbf{326}, 96--192 (2011).

\bibitem{weimer2021} H.~Weimer, A.~Kshetrimayum and R.~Or\'us,
``Simulation methods for open quantum many-body systems,''
\emph{Rev. Mod. Phys.} \textbf{93}, 015008 (2021).

\bibitem{jarkovsky2020} J.~Guth Jarkovsk\'y, A.~Moln\'ar, N.~Schuch and J.~I. Cirac,
``Efficient description of many-body systems with matrix product density operators,''
\emph{PRX Quantum} \textbf{1}, 010304 (2020).

\end{thebibliography}

\end{document}
